% =====================================================================
% Chapitre 1 — Environnement (C2.1.1) et Intégration continue (C2.1.2)
% =====================================================================

\section{Environnement de déploiement et de test}

Le développement de l'application a nécessité un environnement de travail structuré et reproductible. Bien que je sois le seul développeur du projet, je l'ai pensé dès le départ pour faciliter la reprise par n'importe quel développeur~: toute la configuration est versionnée à la racine du dépôt et l'ensemble des services se lance d'une seule commande. L'objectif est double, garantir des conditions d'exécution identiques d'un poste à l'autre et séparer clairement le contexte de développement de celui de la mise en ligne, afin qu'une modification non maîtrisée n'atteigne jamais directement les utilisateurs.

\subsection{Environnement de développement local}\label{subsec:env-local}
Le projet repose sur la stack Angular (frontend), NestJS (backend) et PostgreSQL, avec Prisma comme ORM. Le développement se fait sous Visual Studio Code, complété des extensions adaptées à la stack (Angular Language Service, ESLint, Prettier et Prisma). Les dépendances sont gérées avec npm, le fichier \texttt{package-lock.json} étant versionné pour figer les versions installées d'un poste à l'autre.

Même si le frontend et le backend s'exécutent dans des conteneurs, une version de Node.js cohérente reste requise sur le poste, car l'éditeur et les outils d'analyse statique tournent directement sur l'hôte et les images Docker sont bâties sur une base Node.js. La version est donc figée à Node.js~24 via un fichier \texttt{.nvmrc}, ce qui aligne le poste, les conteneurs et le pipeline d'intégration continue sur une même version.

L'environnement local est orchestré par un unique fichier \texttt{docker-compose.yml} décrivant tous les services interdépendants~: backend, frontend, base PostgreSQL, supervision Grafana et serveur de courriel de test. Tout développeur démarre ainsi le projet complet d'une seule commande, sans rien installer d'autre que Docker. L'option \texttt{develop.watch} reconstruit automatiquement les images quand les fichiers de dépendances changent et synchronise le code à chaud dans les conteneurs pour un rechargement immédiat. La composition de ces services est schématisée en \hyperref[ann:docker-local]{annexe~\ref*{chap:annexe-code}}.

Un second fichier, \texttt{docker-compose.prod.yml}, décrit la variante de production. Il écarte les services de confort du développement et s'appuie sur une image unique, construite en multi-étapes, qui regroupe le frontend compilé servi par nginx et l'API NestJS. Cette approche par conteneurs offre trois avantages~: des conditions d'exécution identiques du poste jusqu'à la mise en ligne, une installation réduite à la seule présence de Docker, et une infrastructure versionnée et traçable au même titre que le code. Le pipeline d'intégration continue exploite ces mêmes définitions en construisant l'image à partir du \texttt{Dockerfile} du dépôt, garantissant que l'image validée par la CI est celle qui sera déployée.

\clearpage
\subsection{Développement, préproduction et production}
Trois contextes d'exécution sont distingués pour séparer les usages et sécuriser le chemin vers la mise en ligne. L'isolation des tests n'est pas un environnement permanent~: elle est assurée par le pipeline d'intégration continue, qui exécute la suite de tests dans des conteneurs éphémères recréés à chaque exécution, sans jamais toucher aux données de développement.

\begin{table}[H]
    \centering
    \begin{tabular}{|>{\raggedright\arraybackslash}p{0.22\linewidth}|>{\raggedright\arraybackslash}p{0.66\linewidth}|}
        \hline
        \textbf{Environnement} & \textbf{Caractéristiques} \\
        \hline
        Développement & Services exécutés localement via Docker Compose, avec rechargement à chaud sur le backend et le frontend. Les variables sensibles (chaîne PostgreSQL, secret \texttt{BETTER\_AUTH\_SECRET}) sont stockées dans des fichiers \texttt{.env} non versionnés. \\
        \hline
        Préproduction & Image de production mise en ligne uniquement pour visualiser le rendu final et illustrer le processus de déploiement. La mise en production réelle étant hors périmètre de ce dossier, cet environnement ne reproduit pas le dispositif cible de l'entreprise. Les migrations Prisma sont appliquées automatiquement au démarrage du conteneur. \\
        \hline
        Production (cible) & Même image déployée via \texttt{docker-compose.prod.yml}, avec des variables injectées sans fichier \texttt{.env} versionné, sur l'hébergeur de l'entreprise (voir 1.1.5). Prévue pour août-septembre 2026. \\
        \hline
    \end{tabular}
\end{table}

\subsection{Outils de qualité et de performance}
Au-delà du lint (ESLint) et du formatage (Prettier), la qualité repose sur un harnais de tests automatisés. Jest couvre les tests unitaires du backend (services et fonctions métier) tandis que Jasmine et Karma, natifs d'Angular, assurent ceux du frontend. Exécutés à la demande en local et systématiquement dans le pipeline d'intégration continue, ces outils garantissent un niveau de qualité constant tout au long du développement.

\subsection{Critères de qualité retenus}
Pour objectiver la notion de qualité, des critères mesurables ont été définis en amont du développement.

\begin{table}[H]
    \centering
    \small
    \begin{tabular}{|>{\raggedright\arraybackslash}p{0.16\linewidth}|>{\raggedright\arraybackslash}p{0.11\linewidth}|>{\raggedright\arraybackslash}p{0.30\linewidth}|>{\raggedright\arraybackslash}p{0.27\linewidth}|}
        \hline
        \textbf{Critère} & \textbf{Cible} & \textbf{Justification} & \textbf{Actuellement} \\
        \hline
        Couverture de tests & 70 \% minimum & Appliquée sur la logique métier la plus sensible du module tournois (calcul du classement, des résultats et des statistiques). & 82,81 \% de statements sur le backend (208 tests), le cœur du module tournois dépassant 90 \%. \\
        \hline
        Erreurs ESLint & 0 en production & Le pipeline d'intégration continue bloque toute fusion vers la branche principale en cas d'erreur de lint. & 0 erreur (backend et frontend). L'application n'étant pas en production, les 11 avertissements frontend restants (labels d'accessibilité, injection) seront corrigés avant la mise en ligne. \\
        \hline
        Temps de réponse API & < 500 ms & Garantir une expérience fluide sur les endpoints principaux. & Jugé satisfaisant lors d'essais manuels (Postman). La mesure automatisée, non encore outillée, sera consolidée avant la production. \\
        \hline
    \end{tabular}
\end{table}

\newpage
\subsection{Choix de l'infrastructure d'hébergement en production}
La plateforme cible est AWS (instance EC2), sur laquelle le \texttt{docker-compose.prod.yml} est déployé via \texttt{docker compose up -d}. Ce choix répond à un objectif de représentativité, AWS étant l'hébergeur utilisé dans mon entreprise, ce qui permet au projet de refléter mon contexte professionnel tout en limitant le risque d'erreur puisque la production reprend un environnement déjà maîtrisé.

Pour ce rendu, la préproduction en ligne n'est toutefois pas hébergée sur AWS. Elle s'appuie sur une solution simple, rapide et peu coûteuse, suffisante pour démontrer le produit et le processus de déploiement sans engager les frais d'une véritable infrastructure de production.

Le critère écoresponsable a par ailleurs fait l'objet d'une analyse comparative, dont il ressort que Scaleway offrirait le meilleur profil.

\begin{table}[H]
    \centering
    \small
    \begin{tabular}{|>{\raggedright\arraybackslash}p{0.20\linewidth}|>{\raggedright\arraybackslash}p{0.22\linewidth}|>{\raggedright\arraybackslash}p{0.20\linewidth}|>{\raggedright\arraybackslash}p{0.20\linewidth}|}
        \hline
        \textbf{Critère} & \textbf{Scaleway (meilleur profil)} & \textbf{Infomaniak} & \textbf{OVHcloud} \\
        \hline
        Énergie & Renouvelable, majoritairement hydraulique & Renouvelable, hydraulique & Renouvelable \\
        \hline
        Refroidissement & Adiabatique, PUE bas & Free cooling, PUE bas & Watercooling, PUE bas \\
        \hline
        Localisation des données & France (UE) & Suisse (hors UE) & France / UE \\
        \hline
        Docker Compose natif & Oui & Oui & Oui \\
        \hline
    \end{tabular}
\end{table}

Scaleway combine énergie renouvelable, refroidissement adiabatique (PUE bas) et hébergement des données en France. Infomaniak présente un profil comparable mais hors Union européenne, et OVHcloud reste une alternative française crédible. Sur ce seul critère, Scaleway s'imposerait. Le choix d'AWS, motivé par la représentativité, n'en demeure pas moins assumé, cette analyse valant recommandation pour une optimisation ultérieure.

Cette architecture présente enfin une limite connue~: packagée dans une image Docker sur une instance unique, elle n'est pas scalable horizontalement. Ce compromis privilégie la rapidité de développement et la simplicité d'installation. Si la charge l'exige, l'évolution prévue en TMA consiste à placer un répartiteur de charge (Elastic Load Balancer) devant plusieurs instances, puis à passer à une orchestration de conteneurs (ECS ou Kubernetes).

\section{Intégration continue}

Afin de détecter les régressions le plus tôt possible dans le cycle de développement, un pipeline d'intégration continue (CI) a été mis en place. Il s'exécute à chaque proposition de fusion vers les branches d'intégration et de production, et empêche l'intégration de code non conforme.

\subsection{Choix de l'outil d'intégration continue}
L'outil retenu est GitLab CI/CD. Ce choix répond d'abord à un objectif de représentativité~: GitLab est la plateforme de stockage des repository utilisée dans mon entreprise, et reproduire cet environnement permet à ce projet de refléter les compétences réellement mises en œuvre dans mon contexte professionnel. Il s'appuie en outre sur une intégration native avec GitLab (le dépôt y étant déjà hébergé), sa gratuité pour les projets publics et la simplicité de configuration des pipelines via un fichier YAML versionné dans le dépôt, ce qui rend la configuration traçable au même titre que le code applicatif.

Si le choix avait été entièrement libre, je me serais orienté vers GitHub Actions~: son écosystème d'actions réutilisables (le Marketplace) accélère nettement la mise en place des pipelines, sa communauté plus large facilite la résolution des problèmes, et son intégration avec de nombreux services tiers en fait l'outil le plus productif pour un projet personnel. Le critère de représentativité a néanmoins primé sur cette préférence.

\subsection{Description du pipeline}\label{subsec:ci-pipeline}
Le pipeline est défini dans le fichier \texttt{.gitlab-ci.yml} et organisé en deux étapes (\emph{stages})~: \texttt{ci} pour la validation du code et \texttt{cd} pour la livraison de l'image. Pour éviter les exécutions redondantes, une règle globale limite son déclenchement aux merge requests ciblant \texttt{main} ou \texttt{dev} et aux étiquettes de version au format \texttt{vX.Y.Z}. Avant chaque job, la version de Node.js est comparée à celle figée dans \texttt{.nvmrc}, afin de garantir un environnement identique à celui du poste et des conteneurs.

L'étape \texttt{ci} exécute trois jobs en parallèle, tandis que l'étape \texttt{cd} n'intervient que lors de la livraison d'une version.

\begin{table}[H]
    \centering
    \small
    \begin{tabular}{|c|>{\raggedright\arraybackslash}p{0.22\linewidth}|>{\raggedright\arraybackslash}p{0.50\linewidth}|}
        \hline
        \textbf{Stage} & \textbf{Job} & \textbf{Rôle} \\
        \hline
        ci & \texttt{ci:backend} & Installation des dépendances (\texttt{npm ci}), génération du client Prisma, lint, build et tests unitaires avec couverture (rapport Cobertura archivé en artefact). \\
        \hline
        ci & \texttt{ci:frontend} & Installation, lint, build et tests Angular en Chromium \emph{headless} avec couverture. \\
        \hline
        ci & \texttt{ci:docker-build} & Construction de l'image Docker de production à partir du \texttt{Dockerfile}, sans publication, pour valider que l'image se construit. \\
        \hline
        cd & \texttt{cd:image} & Sur étiquette \texttt{vX.Y.Z} uniquement~: connexion au registre, construction puis publication de l'image (tags \texttt{vX.Y.Z} et \texttt{latest}). Détaillé au chapitre~4. \\
        \hline
    \end{tabular}
\end{table}

Chaque job est bloquant~: son échec fait échouer le pipeline et empêche la fusion de la merge request concernée. Un extrait commenté du fichier \texttt{.gitlab-ci.yml} et une capture d'une exécution réussie du pipeline sont fournis en \hyperref[ann:ci-extract]{annexe~\ref*{chap:annexe-code}}.

\newpage
\subsection{Stratégie de gestion des branches}
La stratégie de branches structure le flux de contribution et renforce l'efficacité du pipeline.

\begin{table}[H]
    \centering
    \begin{tabular}{|>{\raggedright\arraybackslash}p{0.20\linewidth}|>{\raggedright\arraybackslash}p{0.68\linewidth}|}
        \hline
        \textbf{Branche} & \textbf{Rôle} \\
        \hline
        \texttt{main} & Branche de production, protégée. Aucun push direct~: toute modification passe par une merge request dont le pipeline doit être au vert avant fusion. \\
        \hline
        \texttt{dev} & Branche d'intégration (branche par défaut), protégée, sur laquelle les fonctionnalités sont fusionnées avant d'être portées vers \texttt{main}. \\
        \hline
        \texttt{feat/xxx} & Branches de fonctionnalités, créées depuis \texttt{dev} pour chaque évolution et supprimées une fois fusionnées. \\
        \hline
        \texttt{hotfix/xxx} & Branches de correctif urgent, créées depuis \texttt{main} pour corriger un problème en production, puis fusionnées vers \texttt{main} et reportées sur \texttt{dev}. \\
        \hline
    \end{tabular}
\end{table}

Les branches \texttt{main} et \texttt{dev} sont protégées~: le push direct y est interdit et seule la fusion via merge request est autorisée.

\begin{figure}[H]
    \centering
    \includegraphics[width=0.9\linewidth]{branch}
    \caption{Règles de protection des branches \texttt{main} et \texttt{dev} sur GitLab.}
\end{figure}

La fusion elle-même est conditionnée par des contrôles automatiques~: le pipeline doit être au vert et toutes les discussions doivent être résolues.

\begin{figure}[H]
    \centering
    \includegraphics[width=0.72\linewidth]{mergerequest}
    \caption{Contrôles de fusion imposés avant l'intégration d'une merge request.}
\end{figure}

Cette discipline garantit qu'aucune régression ne peut atteindre la production sans avoir été détectée par le pipeline, et qu'une fonctionnalité en cours n'impacte jamais directement la branche de production. Même seul sur le projet, elle reproduit les conditions d'un travail en équipe et facilite une reprise à plusieurs.
